A Hopfield network with Hebbian weights and the update given in the assignment recalls none of \N{d:n_patterns} dense character bitmaps from noisy copies, and every one of the \N{b:n_queries} queries falls into the same spurious state. This paper reports that baseline as well as a capacity study of (1.) how quickly exact recall collapses as stored pattern similarity grows, (2.) which learning rules hold up, and (3.) why the characters fail. Random patterns that share a prototype, or that have a biased density, are stored with Hebbian, Storkey, covariance, and pseudoinverse rules and with modern (softmax) networks, at loads from 2 to 40 patterns and two noise levels, with 50 pattern sets per condition. The load at which Hebbian exact recall from 10\% noise falls to 50\% (interpolated, 144 neurons) is \N{c:proto:0.0:hebb:10} patterns for independent balanced patterns and \N{c:proto:0.4:hebb:10} at an expected overlap of 0.4. The assignment's update on 0/1 states lowers the first figure to \N{c:proto:0.0:hebb01:10}. Density alone breaks Hebbian storage; the characters also share structure beyond density, and that structure and their uneven densities likely defeat the covariance rule together; and under the assignment's update, 18 patterns are past the 50\% load even for random patterns. The study cannot say how much each of the three contributes. The pseudoinverse rule and a modern network at $\beta = 64$ recall all 18 characters from the test file, but the pseudoinverse recalls \N{f:valid42:pinv:mean} of 18 on average from the 20\% file, over 200 update orders.
